
\section{Objectifs et aire d'étude}

\subsection{Objectifs}


Le projet de recherche a pour but principal d'acquérir des connaissances sur les critères favorisant la connectivité écologique de la petite faune dans la région métropolitaine de Montréal. Les objectifs spécifiques du projet étaient :

\begin{itemize}
    \item de développer des connaissances en écologie routière sur les routes du Ministère ;
    \item d'identifier les critères d'utilisation optimaux des passages fauniques par les rainettes faux-grillons et les tortues (degré d'utilisation, mortalité, connectivité) en fonction des types de routes et des saisons ;
    \item d'améliorer la qualité des habitats des rainettes faux-grillons et des tortues, et d'émettre des recommandations de conception des passages fauniques pour ces groupes (emplacement, protocole, suivi).
\end{itemize}

Afin de répondre à ces objectifs, une première phase de récolte de données s'est déroulée avant les travaux d'élargissement et d'aménagement des passages fauniques en 2024. Durant cette période, l'équipe de recherche a évalué la mortalité routière, ainsi que l'utilisation des routes par la petite faune sur des sections de l'A10 et de la R104, à proximité des emplacements projetés des passages fauniques.

Une fois les travaux d'élargissement et d'aménagement achevés sur ces routes, une deuxième phase du projet serait nécessaire afin d'évaluer l'efficacité des passages fauniques, en quantifiant le degré d'utilisation des passages par la faune selon leurs critères de conception. Puisqu'aucune date de travaux d'élargissement n'est encore prévue, cette deuxième phase a été remplacée par une évaluation des critères favorisant la connectivité écologique des amphibiens, des reptiles et des petits mammifères dans la région métropolitaine de Montréal. Cette partie du projet est présentée à la partie~\ref{part:II} du rapport.

\subsection{Aire d'études}

Le projet de recherche original se déroulait dans la région métropolitaine de Montréal. Les municipalités concernées étaient celles de Brossard et de La Prairie.
La construction de passages fauniques a été planifiée à quatre sites sur l'A10 et à deux sites sur la R104 (Figure 3). L’emplacement des sites a été déterminé par le MTMD, en collaboration avec le MELCCFP et appuyé par les données recueillies en 2022 et 2023 sur les milieux humides et hydriques, les inventaires fauniques, les données hydrologiques, les inventaires des carcasses ou la topographie.